\documentclass[12pt,a4paper]{article}

\usepackage[utf8]{inputenc}
\usepackage[T2A]{fontenc}
\usepackage[russian]{babel}
\usepackage{amsmath,amssymb,amsthm}
\usepackage{array}
\usepackage{booktabs}
\usepackage{xcolor}
\usepackage{geometry}
\usepackage{graphicx}
\usepackage{hyperref}
\usepackage{listings}
\usepackage[most]{tcolorbox}
\usepackage{tikz}
\usepackage{pgfplots}
\usepackage{titlesec}
\pgfplotsset{compat=1.18}

\geometry{left=25mm,right=20mm,top=20mm,bottom=20mm}
\setlength{\parindent}{0pt}
\setlength{\parskip}{0.65em}

\definecolor{TitleBlue}{HTML}{123B66}
\definecolor{SubBlue}{HTML}{2E6DA4}
\definecolor{SoftBlue}{HTML}{EEF5FB}
\definecolor{SoftGreen}{HTML}{F1FAF1}
\definecolor{SoftYellow}{HTML}{FFF9E8}
\definecolor{SoftOrange}{HTML}{FFF3E8}
\definecolor{CodeBg}{HTML}{F8F8F8}
\definecolor{CodeFrame}{HTML}{D0D0D0}

\titleformat{\section}{\Large\bfseries\color{TitleBlue}}{\thesection}{0.9em}{}
\titleformat{\subsection}{\normalfont\large\color{SubBlue}}{\thesubsection}{0.8em}{}
\titleformat{\subsubsection}{\normalfont\bfseries\color{SubBlue!70!black}}{\thesubsubsection}{0.7em}{}

\newtcolorbox{defbox}{
  colback=SoftBlue,
  colframe=TitleBlue!70!black,
  arc=2mm,
  boxrule=0.7pt,
  title=Определение,
  fonttitle=\bfseries
}

\newtcolorbox{formbox}{
  colback=SoftGreen,
  colframe=green!50!black,
  arc=2mm,
  boxrule=0.7pt,
  title=Формула,
  fonttitle=\bfseries
}

\newtcolorbox{notebox}{
  colback=SoftYellow,
  colframe=yellow!55!black,
  arc=2mm,
  boxrule=0.7pt,
  title=Важно,
  fonttitle=\bfseries
}

\newtcolorbox{exbox}{
  colback=SoftOrange,
  colframe=orange!60!black,
  arc=2mm,
  boxrule=0.7pt,
  title=Пример,
  fonttitle=\bfseries
}

\lstdefinestyle{kotline}{
  basicstyle=\ttfamily\small,
  backgroundcolor=\color{CodeBg},
  frame=single,
  rulecolor=\color{CodeFrame},
  numbers=left,
  numberstyle=\tiny\color{gray!80!black},
  stepnumber=1,
  xleftmargin=1.8em,
  xrightmargin=0.6em,
  breaklines=true,
  columns=fullflexible,
  keepspaces=true,
  showstringspaces=false,
  tabsize=2
}

\newtcblisting{codebox}{
  colback=CodeBg,
  colframe=CodeFrame,
  arc=2mm,
  boxrule=0.7pt,
  listing only,
  listing options={style=kotline}
}

\begin{document}

\begin{center}
{\LARGE\bfseries Конспект по интерполяции функций}
\end{center}

\tableofcontents

\section{Введение}
\subsection{Что такое интерполяция}
\begin{defbox}
Интерполяция --- это восстановление значения функции по её значениям в конечном наборе точек. Если заданы узлы $(x_i,y_i)$, где $y_i=f(x_i)$, то требуется построить функцию $P(x)$ такую, что $P(x_i)=y_i$ для всех $i$.
\end{defbox}

Почти всегда в качестве $P(x)$ берут полином или кусочно-полиномиальную функцию. Смысл задачи очень практический: у нас есть дискретные данные, а нам нужно получить значение между ними без прямого вычисления исходной функции.

\subsection{Зачем это нужно}
Интерполяция нужна, когда исходная функция известна только таблично, вычисляется дорого, задаётся измерениями или требуется более гладкая кривая для графика. Ещё одно важное применение --- подготовка численных схем для интегрирования, дифференцирования и решения уравнений.

\subsection{Интерполяция и экстраполяция}
Интерполяция работает внутри диапазона известных узлов. Экстраполяция --- это попытка оценить функцию вне этого диапазона. Экстраполяция обычно менее надёжна, потому что поведение функции за пределами данных может измениться совершенно иначе.

\subsection{Интерполяция и аппроксимация}
Интерполяция обязана проходить точно через все точки. Аппроксимация лишь приближает данные и допускает ошибку в узлах. Это принципиальная разница: при интерполяции мы связываем все точки одной кривой, а при аппроксимации ищем наилучший компромисс.

\subsection{Интерполяционный полином}
Если требуется полином степени не выше $n-1$, проходящий через $n$ различных узлов, то такой полином существует и единственен. Классические методы интерполяции как раз и строят этот полином в разных формах.

\begin{formbox}
Интерполяционный полином --- это не просто красивая кривая, а строго определённая функция, которая совпадает со значениями исходной функции в узлах.
\end{formbox}

\section{Общие термины}
\subsection{Узлы и равномерная сетка}
Узлы --- это точки $x_0,x_1,\dots,x_{n-1}$. Если
\[
 x_{i+1}-x_i=h,
\]
то узлы называются равноотстоящими. Это особенно важно для конечных разностей и центральных формул. При неравномерной сетке приходится использовать более общие схемы, например разделённые разности.

\subsection{Разделённые разности}
\begin{defbox}
Разделённые разности --- это величины, которые позволяют строить интерполяционный полином Ньютона для произвольных узлов.
\end{defbox}

Нулевой и первый порядки определяются так:
\[
f[x_i]=f(x_i)=y_i,
\qquad
f[x_i,x_{i+1}]=\frac{f[x_{i+1}]-f[x_i]}{x_{i+1}-x_i}.
\]
Дальнейшие порядки строятся рекурсивно:
\[
f[x_i,\dots,x_{i+k}] = \frac{f[x_{i+1},\dots,x_{i+k}] - f[x_i,\dots,x_{i+k-1}]}{x_{i+k}-x_i}.
\]

Идея очень похожа на постепенное сведение таблицы к одному коэффициенту. Чем выше порядок разности, тем сильнее она отражает кривизну данных.

\subsection{Конечные разности}
Для равноотстоящих узлов удобнее строить таблицу конечных разностей:
\[
\Delta f_i=f_{i+1}-f_i,\qquad
\Delta^2 f_i=\Delta(\Delta f_i),\qquad
\Delta^k f_i=\Delta(\Delta^{k-1} f_i).
\]
Таблица имеет треугольную форму и является базой для формы Ньютона вперёд, Ньютона назад, а также центральных формул Гаусса, Стирлинга и Бесселя.

\subsection{Параметр $t$}
Если узлы равноотстоящие, вводят безразмерный параметр
\[
 t=\frac{x-x_0}{h}
\]
или центральный вариант относительно среднего узла. Параметр $t$ показывает, насколько далеко точка находится от выбранного узла в единицах шага. Именно через $t$ формулы становятся компактными и легко запоминающимися.

\begin{notebox}
Если сетка неравномерная, конечные разности не подходят. Тогда либо используют разделённые разности, либо переходят к другим методам.
\end{notebox}

\section{Многочлен Лагранжа}
\subsection{Идея, формула и применение}

Метод Лагранжа строит интерполяционный полином через базисные полиномы. Для каждого узла создаётся свой базис $\ell_i(x)$, который равен 1 в собственном узле и 0 во всех остальных. Это очень наглядная идея. Метод применим для любых различных узлов и особенно хорош на небольшом числе точек.

\begin{formbox}
\[
P_{n-1}(x)=\sum_{i=0}^{n-1} y_i\ell_i(x),
\qquad
\ell_i(x)=\prod_{j=0,\,j\ne i}^{n-1}\frac{x-x_j}{x_i-x_j}.
\]
\end{formbox}

Смысл простой: каждый базис отвечает только за одну точку. Когда $x=x_i$, все остальные базисы обращаются в ноль, а $\ell_i(x_i)=1$. Метод работает благодаря разложению по единичным вкладам точек и единственности интерполяционного полинома степени $\le n-1$ через $n$ различных узлов.

\subsection{Ограничения и учебная ценность}

При большом числе узлов полином высокой степени может колебаться, вычисления становятся громоздкими. На практике при больших таблицах предпочитают сплайны или форму Ньютона. Однако для учебных целей Лагранж незаменим: это «каждой точке свой базис». Сначала строятся базисные полиномы, потом они складываются с весами $y_i$.

\subsection{Пошаговый пример с вычислениями}
Возьмём точки $(0,1)$, $(1,2)$, $(2,5)$ и найдём значение при $x=1.5$.

\begin{exbox}
\textbf{Шаг 1. Строим базисные полиномы для $x = 1.5$:}
\[
\ell_0(1.5)=\frac{(1.5-1)(1.5-2)}{(0-1)(0-2)}=\frac{0.5 \cdot (-0.5)}{(-1) \cdot (-2)}=\frac{-0.25}{2}=-0.125
\]
\[
\ell_1(1.5)=\frac{(1.5-0)(1.5-2)}{(1-0)(1-2)}=\frac{1.5 \cdot (-0.5)}{1 \cdot (-1)}=\frac{-0.75}{-1}=0.75
\]
\[
\ell_2(1.5)=\frac{(1.5-0)(1.5-1)}{(2-0)(2-1)}=\frac{1.5 \cdot 0.5}{2 \cdot 1}=\frac{0.75}{2}=0.375
\]

\textbf{Шаг 2. Проверка: $\ell_0 + \ell_1 + \ell_2 = -0.125 + 0.75 + 0.375 = 1.0$ (верно)}

\textbf{Шаг 3. Вычисляем полином:}
\[
P_2(1.5)=y_0 \ell_0(1.5)+y_1 \ell_1(1.5)+y_2 \ell_2(1.5)
\]
\[
P_2(1.5)=1 \cdot (-0.125)+2 \cdot 0.75+5 \cdot 0.375=-0.125+1.5+1.875=3.25
\]

Заметим, что для этого набора точек полином оказался $P_2(x)=x^2+1$. Проверим: $P_2(1.5)=1.5^2+1=2.25+1=3.25$ (верно)
\end{exbox}

\section{Многочлен Ньютона}
\subsection{Идея и формы записи}

Метод Ньютона строит полином последовательно: сначала константа, потом линейный член, затем квадратичный и так далее. Коэффициенты берутся из таблицы разностей. Для произвольных узлов используется форма через разделённые разности:

\begin{formbox}
\[
P_{n-1}(x)=f[x_0]+f[x_0,x_1](x-x_0)+f[x_0,x_1,x_2](x-x_0)(x-x_1)+\dots
\]
\end{formbox}

Преимущество: при добавлении нового узла просто дописывается ещё один член.

\subsection{Форма через конечные разности и выбор направления}

Если узлы равноотстоящие ($x_i=x_0+ih$) и $t=(x-x_0)/h$:
\[
P_{n-1}(x)=f_0+\binom{t}{1}\Delta f_0+\binom{t}{2}\Delta^2 f_0+\dots+\binom{t}{n-1}\Delta^{n-1} f_0,
\]
где $\binom{t}{k}=\frac{t(t-1)\cdots(t-k+1)}{k!}$.

Если точка ближе к началу таблицы, используют формулу вперёд. Если ближе к концу --- формулу назад. Это уменьшает величину множителей и делает вычисления устойчивее. На практике форма Ньютона удобнее Лагранжа: при добавлении новой точки не нужно перестраивать всю сумму, достаточно добавить член и новую строку в таблицу разностей.

\subsection{Ограничение на сетку}

Для конечных разностей нужна строгая равномерность сетки. Если узлы неравномерны, используют разделённые разности.

\subsection{Пошаговый пример с таблицей и вычислениями}

Для точек $(0,1)$, $(1,2)$, $(2,5)$ и $x=1.5$:

\textbf{Таблица конечных разностей:}
\begin{center}
\begin{tabular}{r r r r}
\toprule
$x$ & $f$ & $\Delta f$ & $\Delta^2 f$ \\
\midrule
0 & 1 & 1 & 2 \\
1 & 2 & 3 &  \\
2 & 5 &  &  \\
\bottomrule
\end{tabular}
\end{center}

Параметр $t = \frac{x-x_0}{h} = \frac{1.5-0}{1} = 1.5$. Используем формулу вперёд:
\begin{align*}
P_2(1.5) &= f_0 + \binom{t}{1}\Delta f_0 + \binom{t}{2}\Delta^2 f_0 \\
&= 1 + 1.5 \cdot 1 + \frac{1.5 \cdot 0.5}{2} \cdot 2 \\
&= 1 + 1.5 + \frac{0.75}{2} \cdot 2 \\
&= 1 + 1.5 + 0.75 = 3.25
\end{align*}

Результат совпадает с Лагранжем и с истинным полиномом $P_2(x)=x^2+1$ (верно)

\section{Центральные формулы: Гаусс}
\subsection{Идея и формула}

Формула Гаусса --- это центральная интерполяция для равноотстоящих узлов. Она особенно полезна, когда искомая точка находится рядом с центром таблицы. Используется симметрия относительно середины.

\begin{formbox}
\[
P(x)=f_c+t\overline{\Delta}f_c+\frac{t(t-1)}{2!}\Delta^2 f_{c-1}+\frac{t(t+1)(t-1)}{3!}\Delta^3 f_{c-1}+\dots
\]
где $f_c$ --- центральное значение.
\end{formbox}

Первая формула удобна при $t\ge0$ (точка правее центра), вторая --- при $t<0$ (точка левее центра). Это не две разные теории, а одна центральная идея с разным направлением сдвигов.

\subsection{Когда и почему применять}

Гаусс используют для равноотстоящих узлов в центральной области таблицы. Если точка близка к центру, то формулы, построенные от края, накапливают лишнюю погрешность. Гаусс использует центральную таблицу и работает точнее в середине интервала. Однако если точка сильно смещена к краю, лучше взять Ньютона вперёд или назад. Метод требует нечётного числа точек с явно выраженным центральным узлом.

\section{Центральные формулы: Стирлинг}
\subsection{Идея, формула и применение}

Формула Стирлинга предназначена для интерполяции около центрального узла при нечётном числе узлов. Она использует симметрию таблицы и особенно хороша, когда точка близка к центру.

\begin{formbox}
\[
P(x)=f_0+t\frac{\Delta f_0+\Delta f_{-1}}{2}+\frac{t^2}{2!}\Delta^2 f_{-1}+\frac{t(t^2-1)}{3!}\frac{\Delta^3 f_{-1}+\Delta^3 f_{-2}}{2}+\dots
\]
\end{formbox}

Нечётные разности усредняются, а чётные берутся из центральной области таблицы. Именно поэтому формула выглядит срединной и очень гармоничной.

\subsection{Когда и почему применять}

Стирлинг используют, если узлы равноотстоящие, число узлов нечётное и точка интерполяции находится возле центра таблицы. Если точка находится почти точно в центре, то симметрия данных особенно важна. Формула Стирлинга использует эту симметрию и даёт очень естественное приближение. В центральной области он часто даёт более естественное приближение, чем крайние формулы, потому что использует данные симметрично слева и справа от точки. Если же точка далеко от центра, метод становится менее удобным.

\section{Центральные формулы: Бессель}
\subsection{Идея, формула и применение}

Метод Бесселя применяют, когда точка лежит между двумя центральными узлами (число узлов должно быть чётным). Бессель использует среднее двух центральных значений и полу-сдвиги.

\begin{formbox}
\[
P(x)=\frac{f_{-1}+f_0}{2}+\text{поправки через разности и множители от $t$}.
\]
\end{formbox}

Более точно: метод строится на среднем двух центральных значений и на специальных членах с $t-\frac12$. Это очень удобно, если точка действительно расположена между двумя узлами.

\subsection{Когда и почему применять}

Бессель используют при равноотстоящих узлах, чётном числе точек и положении искомого значения около середины между двумя центральными узлами. Если точка находится между двумя центральными узлами, то симметричная формула должна опираться не на один центр, а на пару центральных значений. Поэтому Бессель использует среднее двух ординат в качестве стартового значения. Он особенно хорош при $t\approx0.5$. Вне центральной области Бессель становится менее удобным. Главный принцип: если Стирлинг работает в центральной точке, то Бессель --- в середине между двумя центральными точками.

\section{Интерполяция сплайнами}
\subsection{Смысл}
Сплайн --- это кусочно-полиномиальная функция. Вместо одного большого полинома строится набор небольших полиномов на каждом отрезке между соседними узлами. Это уже локальная, а не глобальная интерполяция.

\subsection{Кубические сплайны}
Чаще всего используют кубические сплайны. На каждом отрезке $[x_i,x_{i+1}]$ строится кубический полином $S_i(x)$ так, чтобы:
\begin{itemize}
\item $S_i(x_i)=f(x_i)$ и $S_i(x_{i+1})=f(x_{i+1})$;
\item первая производная была непрерывна;
\item вторая производная была непрерывна;
\item на концах были заданы граничные условия, например натуральный сплайн: $S''(x_0)=S''(x_n)=0$.
\end{itemize}

\subsection{Почему это удобно}
Сплайны не дают сильных колебаний на концах, как это бывает у полиномов большой степени. Поэтому для реальных данных это обычно более надёжный выбор.

\subsection{Почему сплайн лучше глобального полинома на практике}
Глобальный полином пытается угодить всем точкам одновременно, и из-за этого может колебаться по всему интервалу. Сплайн локален: если одна точка изменилась, это не рушит всю кривую. Именно поэтому сплайны так популярны в инженерных и вычислительных задачах.

\subsection{Графическая иллюстрация}
\begin{center}
\begin{tikzpicture}[scale=1]
\draw[->] (-0.2,0) -- (5.4,0) node[right] {$x$};
\draw[->] (0,-0.2) -- (0,3.0) node[above] {$y$};
\foreach \x/\y in {0/0.6,1/1.8,2/1.2,3/2.3,4/1.7,5/2.6}
  \fill (\x,\y) circle (1.5pt);
\draw[thick,blue,smooth] plot coordinates {(0,0.6) (1,1.8) (2,1.2) (3,2.3) (4,1.7) (5,2.6)};
\node at (3.2,0.35) {Гладкая кривая через заданные точки};
\end{tikzpicture}
\end{center}

\subsection{Глобальная и локальная интерполяция}
Глобальная интерполяция строит одну функцию на весь интервал. Локальная строит отдельные куски. Сплайны --- это классический пример локальной интерполяции.

\section{Численные примеры}

\subsection{Пример для Лагранжа и Ньютона}

Для точек $(0,1)$, $(1,2)$, $(2,5)$ обе формы дают один и тот же полином $P(x)=x^2+1$. Проверим при $x=1.5$:
\[
P(1.5)=1.5^2+1=2.25+1=3.25.
\]

\subsection{Пример для центральных методов}

Возьмём функцию $f(x)=x^3$ на узлах $-2,-1,0,1,2$ с шагом $h=1$. Значения: $y_0=-8$, $y_1=-1$, $y_2=0$, $y_3=1$, $y_4=8$.

\textbf{Таблица конечных разностей:}
\begin{center}
\begin{tabular}{r r r r r r}
\toprule
$x$ & $f$ & $\Delta f$ & $\Delta^2 f$ & $\Delta^3 f$ & $\Delta^4 f$ \\
\midrule
-2 & -8 & 7 & 6 & 6 & 0 \\
-1 & -1 & 13 & 12 & 6 &  \\
0 & 0 & 25 & 18 &  &  \\
1 & 1 & 43 &  &  &  \\
2 & 8 &  &  &  &  \\
\bottomrule
\end{tabular}
\end{center}

Заметим, что третьи разности равны $6$ (константа), четвёртые разности равны $0$. Это правильно, потому что исходная функция --- кубический полином: $\Delta^3(x^3) = 6$, $\Delta^4(x^3) = 0$.

\textbf{Применим метод Стирлинга для $x = -0.3$ (близко к центру $x=0$):}

Центр таблицы: $x_c = 0$, $y_c = 0$. Параметр $t = \frac{-0.3 - 0}{1} = -0.3$.

Начисляем первое приближение:
\[
\overline{\Delta} f_0 = \frac{\Delta f_0 + \Delta f_{-1}}{2} = \frac{25 + 13}{2} = 19
\]
\[
P(-0.3) \approx 0 + (-0.3) \cdot 19 = -5.7
\]

Второй порядок (чётная разность):
\[
P(-0.3) \approx -5.7 + \frac{(-0.3)^2}{2} \cdot 12 = -5.7 + 0.045 \cdot 12 = -5.7 + 0.54 = -5.16
\]

Истинное значение: $f(-0.3) = (-0.3)^3 = -0.027$. Это грубое приближение, но оно показывает, как метод работает.

\subsection{Пример сплайна}

Для точек $(0,0)$, $(1,1)$, $(2,0)$ построим натуральный кубический сплайн. На интервале $[0,1]$ имеем 1 кубический полином, на $[1,2]$ --- другой. Условия:
\begin{itemize}
\item На $x=0$ и $x=2$ вторые производные равны нулю (натуральный сплайн).
\item Непрерывность первой и второй производных при $x=1$.
\item Прохождение через все три точки.
\end{itemize}

Решение этой системы даёт два кубических полинома $S_0(x)$ и $S_1(x)$, которые вместе образуют гладкую кривую через точки.

\section{Сравнение методов}
\begin{center}
\begin{tabular}{p{3.2cm} p{4.0cm} p{5.1cm}}
\toprule
Метод & Когда применять & Плюсы и минусы \\
\midrule
Лагранж & Любые различные узлы, небольшое число точек & Простая явная формула; при большом числе узлов становится громоздким и может быть неустойчивым. \\
Ньютон & Табличные данные, равноотстоящие узлы или произвольные узлы через разделённые разности & Легко добавлять новые узлы; удобно для вычислений; нужно правильно выбрать forward/backward или разделённые разности. \\
Гаусс & Равноотстоящие узлы, центральная область таблицы & Хорошая точность возле центра; нужно правильно выбрать первую или вторую формулу. \\
Стирлинг & Равноотстоящие узлы, нечётное число точек, точка почти в центре & Симметричен и удобен в центре; хуже работает, если точка далеко от центра. \\
Бессель & Равноотстоящие узлы, чётное число точек, точка между двумя центральными узлами & Очень удобен при $t\approx0.5$; менее удобен вне центральной области. \\
Сплайны & Практические данные, когда нужна гладкая локальная кривая & Устойчивы и локальны; требуют решения СЛАУ, но обычно это лучший практический выбор. \\
\bottomrule
\end{tabular}
\end{center}

\subsection{Как выбрать метод по ситуации}
\begin{itemize}
\item Нужна простая учебная формула --- Лагранж.
\item Есть таблица с равным шагом --- Ньютон.
\item Точка около центра --- Гаусс или Стирлинг.
\item Точка между двумя центральными узлами --- Бессель.
\item Важна устойчивость на реальных данных --- сплайны.
\end{itemize}

\section{Ответы на 17 контрольных вопросов}
\begin{enumerate}
\item \textbf{Когда возникает необходимость в использовании интерполяционных методов?}

Интерполяция нужна, когда функция известна только в нескольких точках и нужно восстановить её значение между ними. Это полезно для таблиц, измерений, графиков и численных методов.

\item \textbf{Чем отличается аппроксимация от интерполяции?}

Интерполяция проходит точно через все точки, а аппроксимация лишь приближает их с минимальной ошибкой по выбранному критерию.

\item \textbf{В чём сущность задачи интерполирования?}

Суть --- построить функцию, которая совпадает со значениями в узлах и даёт разумное поведение между ними.

\item \textbf{Поясните смысл терминов: интерполяция, экстраполяция.}

Интерполяция --- внутри известного диапазона, экстраполяция --- вне него. Экстраполяция опаснее и обычно менее точна.

\item \textbf{Как найти приближенное значение функции при линейной интерполяции?}

Используют формулу прямой через две точки:
\[
P(x)=y_0+\frac{y_1-y_0}{x_1-x_0}(x-x_0).
\]

\item \textbf{Как найти приближенное значение функции при квадратичной интерполяции?}

Берут три точки и строят полином второй степени, например по Лагранжу:
\[
P_2(x)=y_0\ell_0(x)+y_1\ell_1(x)+y_2\ell_2(x).
\]

\item \textbf{Как строится интерполяционный многочлен Лагранжа?}

Строятся базисные полиномы $\ell_i(x)$, которые равны 1 в собственном узле и 0 во всех остальных. Затем их линейная комбинация с весами $y_i$ даёт итоговый полином.

\item \textbf{Дайте определение понятий разделенной разности нулевого и первого порядков.}

Нулевой порядок --- это значение функции в узле, первый порядок --- отношение разности значений к разности аргументов.

\item \textbf{Объясните принцип построения интерполяционного полинома Ньютона.}

Полином строится по степеням вида $(x-x_0)$, $(x-x_0)(x-x_1)$ и так далее, а коэффициенты берутся из таблицы разделённых или конечных разностей.

\item \textbf{Покажите графическую интерпретацию интерполяции.}

Интерполяция --- это кривая через заданные точки. Для сплайнов кривая строится по частям, но остаётся гладкой.

\item \textbf{В каких случаях используются конечные разности, в каких – разделённые?}

Конечные разности --- для равноотстоящих узлов. Разделённые разности --- для произвольных узлов.

\item \textbf{В каких случаях используют формулу Ньютона для интерполирования вперед и для интерполирования назад?}

Вперёд --- когда точка ближе к началу таблицы, назад --- когда точка ближе к концу.

\item \textbf{В каких случаях используют формулу Гаусса для интерполирования вперед и для интерполирования назад?}

Гаусс применяют около центра таблицы. Первая формула удобна при $t\ge0$, вторая --- при $t<0$.

\item \textbf{В каких случаях используют формулу Стирлинга?}

Стирлинг используют при нечётном числе равноотстоящих узлов и точке, близкой к центру таблицы.

\item \textbf{В каких случаях используют формулу Бесселя?}

Бессель применяют при чётном числе равноотстоящих узлов, когда точка лежит между двумя центральными узлами.

\item \textbf{В чем разница между глобальной и локальной разновидностями интерполяции?}

Глобальная интерполяция строит одну функцию на весь интервал. Локальная строит куски отдельно, как сплайны. Локальные методы обычно устойчивее.

\item \textbf{В чем заключается интерполяция кубическими сплайнами?}

Кубический сплайн --- это кусочно-кубическая функция, которая проходит через все узлы, имеет непрерывные первую и вторую производные и обычно задаётся натуральными или другими граничными условиями.
\end{enumerate}

\section{Объяснение кода}

\subsection{\texttt{LagranzhSolver.kt} --- реализация метода Лагранжа}

\begin{codebox}
override fun solve(request: InterpolationRequestDto): InterpolationResult {
    val points = normalizePoints(request.points)
    val n = points.size

    val evaluate: (BigDecimal) -> BigDecimal = { calcX ->
        var res = BigDecimal.ZERO
        for (i in 0 until n) {
            var li = BigDecimal.ONE
            for (j in 0 until n) {
                if (j != i) {
                    val numerator = calcX - points[j].x
                    val denominator = points[i].x - points[j].x
                    li = (li * numerator).divide(denominator, mc)
                }
            }
            res += points[i].y * li
        }
        res
    }
    val result = evaluate(request.x)
    return InterpolationResult(type = InterpolationType.LAGRANZH, x = request.x.round(), y = result.round())
}
\end{codebox}

Логика реализации прямо соответствует математической формуле. Функция \texttt{evaluate} реализует полином Лагранжа для произвольного $x$. Нормализация входных точек: сортировка по $x$ и проверка, что абсциссы различны. Внешний цикл по $i$ итерирует по каждому базисному полиному $\ell_i(x)$. Внутренний цикл по $j$ строит произведение множителей $(x-x_j)/(x_i-x_j)$. Используется \texttt{BigDecimal} для контроля округления при делении. Результат --- сумма всех взвешенных базисов с весами $y_i$.

\subsection{\texttt{NewtonSolver.kt} --- последовательное построение полинома}

\begin{codebox}
override fun solve(request: InterpolationRequestDto): InterpolationResult {
    val points = normalizePoints(request.points)
    val x = request.x
    val n = points.size

    checkEquidistant(points)
    val h = points[1].x - points[0].x
    val diffs = InterpolationSolver.calculateFiniteDifferences(points)

    val evaluate: (BigDecimal) -> BigDecimal = { calcX ->
        val useFwd = (calcX - points[0].x).abs() <= (calcX - points.last().x).abs()
        val tCalc = if (useFwd) (calcX - points[0].x).divide(h, mc) 
                    else (calcX - points.last().x).divide(h, mc)
        var res = if (useFwd) points[0].y else points.last().y
        var termCalc = BigDecimal.ONE
        for (k in 1 until n) {
            val delta = if (useFwd) tCalc - BigDecimal(k - 1) 
                       else tCalc + BigDecimal(k - 1)
            termCalc = termCalc.multiply(delta, mc)
            val deltaY = if (useFwd) diffs[k][0] else diffs[k][n - 1 - k]
            res = res.add(termCalc.multiply(deltaY, mc).divide(factorial(k, mc), mc))
        }
        res
    }
    val result = evaluate(x)
    return InterpolationResult(type = InterpolationType.NEWTON, x = x, y = result)
}
\end{codebox}

Проверка равномерности: \texttt{checkEquidistant()} требует одинакового шага $h$. Выбор направления: если точка ближе к левому краю, используется формула вперёд (вычисляем от $x_0$), если ближе к правому --- назад. Таблица разностей вычисляется один раз, затем используется в цикле. Параметр $t$ нормируется: $t = (x - x_0) / h$ или $t = (x - x_n) / h$. Полином строится последовательно: стартовое значение плюс поправки через произведения множителей. Каждый множитель содержит множители типа $(t - (k-1))$ для формулы вперёд или $(t + (k-1))$ для назад. Разности берутся из подготовленной таблицы.

\subsection{\texttt{GaussSolver.kt} --- центральная интерполяция}

Гаусс требует нечётного числа узлов. Центральный индекс вычисляется как \texttt{centerIdx = n / 2}. Выбор между первой и второй формулой зависит от знака параметра $t = (x - x_0) / h$ относительно центра. Разности выбираются со смещениями (\texttt{offset}) в зависимости от порядка и выбранной формулы. Для первой формулы при $t \geq 0$ разности берутся справа от центра, для второй при $t < 0$ --- слева.

\subsection{\texttt{StirlingSolver.kt} --- симметричная центральная формула}

Стирлинг применяется при нечётном числе узлов. Ключевая особенность: нечётные разности усредняются. Начало полинома содержит центральное значение и линейный член, использующий среднюю первую разность слева и справа: \texttt{(diffs[1][centerIdx-1] + diffs[1][centerIdx]) / 2}. Чётные порядки берут центральные разности напрямую. Нечётные порядки усредняются, что обеспечивает симметрию метода и лучшую точность в центральной области.

\subsection{\texttt{BesselSolver.kt} --- интерполяция между двумя центрами}

Бессель использует два центральных значения, которые усредняются как начальное приближение. Число узлов должно быть чётным. Начало: \texttt{result = (y0 + y1) / 2}. Характерный множитель для первой разности содержит полу-сдвиг: \texttt{t - 0.5}. Для чётных порядков разности также усредняются (берутся соседние). Для нечётных порядков разности берутся без усреднения. Метод особенно эффективен, когда точка интерполяции находится почти ровно между двумя центральными узлами (т.е. $t \approx 0.5$).

\subsection{Общий интерфейс и вспомогательные методы}

Все солверы наследуют интерфейс \texttt{InterpolationSolver}, который определяет \texttt{solve()}, \texttt{supports()}, и вспомогательные методы:
\begin{itemize}
\item \texttt{normalizePoints()} --- сортирует узлы по $x$ и проверяет их различимость.
\item \texttt{checkEquidistant()} --- проверяет равномерный шаг для методов, которые это требуют.
\item \texttt{factorial(k)} --- вычисляет $k!$ через \texttt{BigDecimal}.
\item \texttt{calculateFiniteDifferences()} --- строит треугольную таблицу конечных разностей $\Delta^j f_i$.
\item \texttt{generateGraphicPoints()} --- создаёт точки для рисования кривой полинома между узлами.
\end{itemize}

Все методы используют \texttt{BigDecimal} с \texttt{MathContext(16, RoundingMode.HALF\_UP)} для контроля точности на 16 значащих цифр. Это критично: одна маленькая ошибка в промежуточных произведениях может разрастись в старших членах полинома.

\section{Схематичные иллюстрации}
\subsection{Интерполяция через точки}
\begin{center}
\begin{tikzpicture}[scale=1]
\draw[->] (-0.2,0) -- (5.4,0) node[right] {$x$};
\draw[->] (0,-0.2) -- (0,3.0) node[above] {$y$};
\foreach \x/\y in {0/0.6,1/1.8,2/1.2,3/2.3,4/1.7,5/2.6}
  \fill (\x,\y) circle (1.5pt);
\draw[thick,blue,smooth] plot coordinates {(0,0.6) (1,1.8) (2,1.2) (3,2.3) (4,1.7) (5,2.6)};
\end{tikzpicture}
\end{center}

\subsection{Идея локальной интерполяции}
\begin{center}
\begin{tikzpicture}[scale=1]
\draw[->] (-0.2,0) -- (5.4,0) node[right] {$x$};
\draw[->] (0,-0.2) -- (0,3.0) node[above] {$y$};
\foreach \x/\y in {0/0.6,1/1.8,2/1.2,3/2.3,4/1.7,5/2.6}
  \fill (\x,\y) circle (1.5pt);
\draw[thick,red] (0,0.6) .. controls (0.6,1.0) and (0.6,1.8) .. (1,1.8)
                 .. controls (1.5,1.7) and (1.5,1.1) .. (2,1.2)
                 .. controls (2.5,1.4) and (2.5,2.2) .. (3,2.3)
                 .. controls (3.5,2.4) and (3.5,1.7) .. (4,1.7)
                 .. controls (4.5,1.8) and (4.5,2.5) .. (5,2.6);
\end{tikzpicture}
\end{center}

\section*{Заключение}
В этом файле собраны базовые теоретические разделы по интерполяции функций, подробные разборы методов Лагранжа, Ньютона, Гаусса, Стирлинга, Бесселя и сплайнов, сравнительная таблица, ответы на контрольные вопросы и пояснение кода. Документ можно открывать в Overleaf и дорабатывать при необходимости.

\end{document}
